\chapter{Evaluation Metrics}
\label{app:ml_metrics}

This appendix defines the binary-classification metrics used to evaluate the Isolation Forest model in Chapter 5. For each formula, TP, FP, FN and TN represent the four cells of the confusion matrix (Table~\ref{tab:confusion}): True Positive, False Positive, False
Negative, and True Negative respectively.

% ─────────────────────────────────────────────────────────────────────────────
\section{Confusion Matrix Cells}

Given a binary classifier that labels each sample as either
\textit{Anomaly} (positive class) or \textit{Normal} (negative class),
the four symbols used throughout Chapter~5 are defined in
Table~\ref{tab:cm_defs}:

\begin{table}[H]
\centering
\caption{Confusion matrix cell definitions}
\label{tab:cm_defs}
\small
\begin{tabular}{llp{8.5cm}}
\toprule
\textbf{Symbol} & \textbf{Name} & \textbf{Definition} \\
\midrule
TP & True Positive  & Anomaly sample correctly predicted as Anomaly \\
FP & False Positive & Normal sample incorrectly predicted as Anomaly \\
FN & False Negative & Anomaly sample incorrectly predicted as Normal \\
TN & True Negative  & Normal sample correctly predicted as Normal \\
\bottomrule
\end{tabular}
\end{table}

% ─────────────────────────────────────────────────────────────────────────────
\section{Classification Metrics}

\subsection*{Precision}

Precision measures what fraction of the samples predicted as anomalous
are actually anomalous. High precision means few false alarms.

\begin{equation}
\label{eq:precision}
  \text{Precision} = \frac{\text{TP}}{\text{TP} + \text{FP}}
\end{equation}

\textit{From the confusion matrix (Table~\ref{tab:confusion}, benchmark HR-only model):} Anomaly precision $\approx 367/(367+14) \approx 0.96$.

\subsection*{Recall (Sensitivity)}

Recall measures the fraction of all true anomalies recovered. In a health-monitoring context, a missed cardiac event (false negative) carries more risk than a spurious alert, making recall the more safety-sensitive of the two metrics.

\begin{equation}
\label{eq:recall}
  \text{Recall} = \frac{\text{TP}}{\text{TP} + \text{FN}}
\end{equation}

\textit{From the confusion matrix (Table~\ref{tab:confusion}, benchmark HR-only model):} Anomaly recall $\approx 367/(367+159) \approx 0.70$.

\subsection*{F1-Score}

F1 is the harmonic mean of precision and recall, and is the standard summary metric when the class distribution is skewed --- here, 100 healthy vs.\ 526 diseased samples in the UCI held-out test set.

\begin{equation}
\label{eq:f1}
  \text{F1} = 2 \cdot \frac{\text{Precision} \times \text{Recall}}
                             {\text{Precision} + \text{Recall}}
            = \frac{2\,\text{TP}}{2\,\text{TP} + \text{FP} + \text{FN}}
\end{equation}

\textit{From the confusion matrix (Table~\ref{tab:confusion}, benchmark HR-only model):} Anomaly F1 $= 2 \times 0.96 \times 0.70 /
(0.96 + 0.70) \approx 0.81$.

\subsection*{Weighted Average}

The weighted average reported in Table~\ref{tab:ml_eval} computes each metric as the
support-weighted mean across both classes:

\begin{equation}
\label{eq:weighted_avg}
  M_{\text{weighted}} =
    \frac{n_{\text{Normal}} \cdot M_{\text{Normal}}
        + n_{\text{Anomaly}} \cdot M_{\text{Anomaly}}}
         {n_{\text{Normal}} + n_{\text{Anomaly}}}
\end{equation}

where $M$ is the metric of interest (Precision, Recall, or F1) and
$n_c$ is the support (number of samples) for class $c$.

% ─────────────────────────────────────────────────────────────────────────────
\section{Contamination Hyperparameter}

The \texttt{contamination} parameter $\rho \in (0, 0.5)$ passed to
scikit-learn's \texttt{IsolationForest} sets the decision threshold
$\tau$ as the $\rho$-th percentile of anomaly scores on the training
set~\cite{Pedregosa2011}:

\begin{equation}
\label{eq:contamination}
  \tau = \text{Percentile}_{\rho}
         \bigl(\{\,s_i \mid x_i \in X_{\text{train}}\,\}\bigr)
\end{equation}

Given a sample $x$, its score is $s = \texttt{decision\_function}(x)$. If $s < \tau$, then $x$ is classified as an anomaly. Increasing $\rho$ lowers $\tau$, which increases recall at the cost of more false positives. The production value $\rho = 0.020$ (2\%) is chosen via grid search on the combined UCI and synthetic training data (Section~\ref{sec:anomaly_design}), prioritising a low false-positive rate for continuous home monitoring. It is independent of the UCI held-out benchmark (Section~\ref{sec:ml_validation}), which uses a separate configuration ($\rho = 0.10$, $n_{\text{estimators}} = 200$).


